\section{Why LLMs Work: A Theoretical Foundation}
\label{sec:why-llms-work}

There is currently no widely accepted mechanistic theory of why next-token
prediction on text produces general intelligence. The standard story is
empirical: scaling laws hold, capabilities emerge, benchmarks improve.
Other theoretical programs exist --- information-theoretic accounts,
compression views, in-context learning theory, grokking --- but none has
achieved consensus as a mechanistic explanation of what the model is
computing.

\subsection*{The BP Interpretation}

If transformers implement belief propagation over an implicit factor graph,
then next-token prediction can be interpreted as marginal inference in a
Bayesian network over language concepts. Under this interpretation, the
model behaves as though it is learning a structured probabilistic model of
the world and performing inference in it.

This suggests a mechanistic account of what the model is computing: each
forward pass propagates beliefs through an implicit factor graph, updating
each proposition's probability given the evidence in the context window.
Next-token prediction is the output of this inference process --- the most
probable next token given the current belief state.

\subsection*{Scaling Through the BP Lens}

The scaling dimensions have natural BP interpretations:

\begin{center}
\begin{tabular}{ll}
\toprule
Dimension & BP Interpretation \\
\midrule
$W$ (sequence length) & Nodes in the factor graph \\
$L$ (layers) & Rounds of belief propagation \\
$H$ (attention heads) & AND patterns per round \\
$D_{ff}$ (FFN hidden units) & OR patterns per round \\
$D_{\mathrm{model}}$ (embedding) & Belief state capacity \\
\bottomrule
\end{tabular}
\end{center}

Scaling may improve performance because it expands inference capacity ---
more rounds, richer potentials, more simultaneous concepts. This is
qualitatively consistent with the BP interpretation, though the specific
power-law exponents are not derived from BP theory.

\subsection*{What This Does Not Explain}

The BP framework suggests a structure for what the model computes. It does
not explain why the training distribution (internet text) produces a useful
factor graph, nor does it derive the power-law scaling exponents from first
principles. These remain open questions.

\subsection*{Epistemic Status}

\begin{center}
\begin{tabular}{ll}
\toprule
Claim & Status \\
\midrule
Sigmoid transformers can implement BP & Formal theorem \\
Trained LLMs implement BP & Empirical hypothesis \\
Next-token prediction fits a Bayes net & Mechanistic interpretation \\
Scaling expands BP inference capacity & Mechanistic interpretation \\
BP explains power-law scaling exponents & Not established \\
\bottomrule
\end{tabular}
\end{center}